\section{Finite-flat periodic schemes}

Let $B=\Z[A,A^{-1}]$ and let $H_A$ denote \eqref{eq:henon} over $B$.

\begin{theorem}[Cyclic presentation and rank]\label{thm:flat}
For every $n\geq1$, the fixed scheme $\mathcal X_n=\Fix(H_A^n)$ is
scheme-theoretically isomorphic to
\begin{equation}\label{eq:cyclicscheme}
 \Spec B[x_0,\ldots,x_{n-1}]/(f_0,\ldots,f_{n-1}),
 \qquad
 f_i=A x_i^2-1+\sum_{j\in\{i-1,i+1\}}x_j,
\end{equation}
where neighbors are counted as a multiset modulo $n$.  The morphism
$\mathcal X_n\to\Spec B$ is finite flat of rank $2^n$.
\end{theorem}

\begin{proof}
Writing $H_A^i(x_0,x_{-1})=(x_i,x_{i-1})$ gives the recurrence in
\eqref{eq:henon}.  The scheme map sends $(q,p)$ to
$x_i=\pi_1H_A^i(q,p)$; its inverse sends a cyclic tuple to
$(x_0,x_{n-1})$.  The recurrence verifies both compositions on coordinate
rings.  For $n=1$ the inverse is $x_0\mapsto(x_0,x_0)$, and for $n=2$ it is
$(x_0,x_1)\mapsto(x_0,x_1)$.  In particular, the multiset convention gives
\[
 n=1:\quad Ax_0^2+2x_0-1,
\]
and
\[
 n=2:\quad Ax_0^2+2x_1-1,\qquad Ax_1^2+2x_0-1.
\]
After division by the unit $A$, every $f_i$ is monic with leading monomial
$x_i^2$ under any degree-compatible order.  These monomials are pairwise
coprime, so Buchberger's product criterion applies over $B$.  The standard
monomials are $\prod_i x_i^{e_i}$ with $e_i\in\{0,1\}$, giving a free
$B$-basis of size $2^n$.
\end{proof}

The theorem separates total scheme length from geometric support size.  The
latter equals $2^n$ only at a reduced fiber.

\begin{proposition}[Degree-drop primes]\label{prop:degree-drop}
For an integer $a\ne0$, every prime $p\nmid a$ has fiber scheme length $2^n$.
A prime $p\mid a$ is not a point of
$\Spec\Z[A,A^{-1}]$.  Direct reduction of the original uninverted
$\Z[A]$-family instead gives
\[
 H_0(q,p)=(1-p,q),\qquad H_0^4=I.
\]
In particular, $\Fix(H_0^n)=\A^2$ whenever $4\mid n$.
\end{proposition}

Thus `degree-good' means $p\nmid a$; it does not mean \`etale.  Let $J_n$ be
the Jacobian of the cyclic equations and let $M_n$ be the derivative monodromy
of the H\'enon orbit.  Linearizing the scalar recurrence identifies
$\ker J_n$ with $\ker(I-M_n)$: a cyclic tangent vector satisfies
\[
 \binom{v_{i+1}}{v_i}
 =DH_a(x_i,x_{i-1})\binom{v_i}{v_{i-1}}.
\]
Thus a degree-good fiber is \`etale precisely when no geometric periodic
point has multiplier one.  Every degree-good characteristic-two fiber
($a\ne0$) is non-\`etale, since then
\[
 DH_a=\begin{pmatrix}0&1\\1&0\end{pmatrix}
\]
and every power has eigenvalue one.  The qualifier is necessary: at $a=0$
some fixed schemes are empty, hence \`etale, whereas $\Fix(H_0^4)=\A^2$.

For later use, exact norm calculations at the first two periods give
\begin{equation}\label{eq:discriminants}
 D_{a,1}=-4(a+1),\qquad
 D_{a,2}=2^8(a+1)(a-3)^3.
\end{equation}
